% GENERATED FILE. Edit canonical lessons, then run npm run generate:books.
% canonical-source-hash: fnv1a64:1532a919ed3dfbe2
% canonical-lessons: TA-C400-nataimetai, TA-C400-tantavalam, TA-C400-varukai, TA-C400-ottivai, TA-C400-maru-change

\chapter{Platforms and Tracks}
\label{ch:ta-platforms-and-tracks}
\hlchaptermodality{400}

\begin{chapteropening}
\textbf{By the end of this chapter:} \emph{I can talk about a platform, a railway track, an arrival and a trip put off.}

Complete the last lesson of chapter 400: I can talk about a platform, a railway track, an arrival and a trip put off.
\end{chapteropening}

\section[\texorpdfstring{\ta{நடைமேடை} (naṭaimēṭai)}{நடைமேடை (naṭaimēṭai)}]{\ta{நடைமேடை} (naṭaimēṭai) — a platform; a raised walkway}

\label{lesson:TA-C400-nataimetai}



\begin{quote}
Before the new one: say the Tamil for a complaint, then the Tamil for an athlete.
\end{quote}

\subsection*{You'll want to know: \ta{நடைமேடை}}
\textbf{\ta{நடைமேடை}} — \emph{naṭaimēṭai} — \textquotedblleft{}a platform; a raised walkway\textquotedblright{}.

\begin{quote}
\ta{ரயில்} \ta{மூன்றாவது} \ta{நடைமேடைக்கு} \ta{வருகிறது}.
\end{quote}

\begin{quote}
\emph{rayil mūṉṟāvatu naṭaimēṭaikku varukiṟatu.}
\end{quote}

\begin{quote}
The train is coming to the third platform.
\end{quote}

Before -\ta{க்கு} the final -\ta{ஐ} stays as it is: \ta{நடைமேடைக்கு}, to the platform.

\begin{cousinweb}
\textbf{\ta{நடைமேடை}} (\emph{naṭaimēṭai}) is built from \emph{naṭai}, \textquotedblleft{}walking, gait\textquotedblright{} (DEDR 3582, the family of \ta{நட}, to walk), and \emph{mēṭai}, \textquotedblleft{}a raised floor, a stage\textquotedblright{}, whose origin is not settled.
\end{cousinweb}

\begin{grammarlens}[title={What you've built}]
One more word for getting around.
\end{grammarlens}

\subsection*{Your turn}
\begin{itemize}
\raggedright
  \item \emph{Say it:} \emph{naṭaimēṭai}
  \item \emph{Say it:} \emph{naṭaimēṭai}, once more
  \item \emph{Say it:} \emph{rayil mūṉṟāvatu naṭaimēṭaikku varukiṟatu.} — The train is coming to the third platform.
  \item \emph{Recall:} say the Tamil for a complaint, then the Tamil for an athlete, then say \emph{naṭaimēṭai} again
\end{itemize}

\begin{tcolorbox}[breakable,colback=teal!4,colframe=teal!35!black,title={Before you move on}]
What does \ta{நடைமேடை} mean? (\textquotedblleft{}A platform; a raised walkway\textquotedblright{}.) Say it once more.
\end{tcolorbox}
\section[\texorpdfstring{\ta{தண்டவாளம்} (taṇṭavāḷam)}{தண்டவாளம் (taṇṭavāḷam)}]{\ta{தண்டவாளம்} (taṇṭavāḷam) — a railway track; a rail}

\label{lesson:TA-C400-tantavalam}



\begin{quote}
Before the new one: say the Tamil for an athlete, then the Tamil for a platform.
\end{quote}

\subsection*{You'll want to know: \ta{தண்டவாளம்}}
\textbf{\ta{தண்டவாளம்}} — \emph{taṇṭavāḷam} — \textquotedblleft{}a railway track; a rail\textquotedblright{}.

\begin{quote}
\ta{குழந்தைகள்} \ta{தண்டவாளத்தில்} \ta{நடக்கக்} \ta{கூடாது}.
\end{quote}

\begin{quote}
\emph{kuḻantaikaḷ taṇṭavāḷattil naṭakkak kūṭātu.}
\end{quote}

\begin{quote}
Children must not walk on the railway track.
\end{quote}

A -\ta{அம்} noun takes -\ta{அத்து}- before a case ending: \ta{தண்டவாளத்தில்}, on the rails.

\begin{cousinweb}
\textbf{\ta{தண்டவாளம்}} (\emph{taṇṭavāḷam}) is listed in DEDR 3050 with the senses \textquotedblleft{}cast iron, an iron rail, a girder\textquotedblright{}; where the word comes from before that is not settled.
\end{cousinweb}

\begin{grammarlens}[title={What you've built}]
One more word for getting around.
\end{grammarlens}

\subsection*{Your turn}
\begin{itemize}
\raggedright
  \item \emph{Say it:} \emph{taṇṭavāḷam}
  \item \emph{Say it:} \emph{taṇṭavāḷam}, once more
  \item \emph{Say it:} \emph{kuḻantaikaḷ taṇṭavāḷattil naṭakkak kūṭātu.} — Children must not walk on the railway track.
  \item \emph{Recall:} say the Tamil for an athlete, then the Tamil for a platform, then say \emph{taṇṭavāḷam} again
\end{itemize}

\begin{tcolorbox}[breakable,colback=teal!4,colframe=teal!35!black,title={Before you move on}]
What does \ta{தண்டவாளம்} mean? (\textquotedblleft{}A railway track; a rail\textquotedblright{}.) Say it once more.
\end{tcolorbox}
\section[\texorpdfstring{\ta{வருகை} (varukai)}{வருகை (varukai)}]{\ta{வருகை} (varukai) — arrival; a visit}

\label{lesson:TA-C400-varukai}



\begin{quote}
Before the new one: say the Tamil for a platform, then the Tamil for a railway track.
\end{quote}

\subsection*{You'll want to know: \ta{வருகை}}
\textbf{\ta{வருகை}} — \emph{varukai} — \textquotedblleft{}arrival; a visit\textquotedblright{}.

\begin{quote}
\ta{விமானத்தின்} \ta{வருகை} \ta{ஒரு} \ta{மணி} \ta{நேரம்} \ta{தாமதமானது}.
\end{quote}

\begin{quote}
\emph{vimāṉattiṉ varukai oru maṇi nēram tāmatamāṉatu.}
\end{quote}

\begin{quote}
The plane's arrival was an hour late.
\end{quote}

\ta{வருகை} is a noun made from \ta{வா}, to come: its stem \ta{வரு}- takes the ending -\ta{கை}.

\begin{cousinweb}
\textbf{\ta{வருகை}} (\emph{varukai}) is built on \emph{varu-}, the stem of \ta{வா}, \textquotedblleft{}to come\textquotedblright{}, with the noun ending \emph{-kai}; DEDR 5270 gathers the family, including \ta{வரவு}, \textquotedblleft{}coming, income\textquotedblright{}.
\end{cousinweb}

\begin{grammarlens}[title={What you've built}]
One more word for getting around.
\end{grammarlens}

\subsection*{Your turn}
\begin{itemize}
\raggedright
  \item \emph{Say it:} \emph{varukai}
  \item \emph{Say it:} \emph{varukai}, once more
  \item \emph{Say it:} \emph{vimāṉattiṉ varukai oru maṇi nēram tāmatamāṉatu.} — The plane's arrival was an hour late.
  \item \emph{Recall:} say the Tamil for a platform, then the Tamil for a railway track, then say \emph{varukai} again
\end{itemize}

\begin{tcolorbox}[breakable,colback=teal!4,colframe=teal!35!black,title={Before you move on}]
What does \ta{வருகை} mean? (\textquotedblleft{}Arrival; a visit\textquotedblright{}.) Say it once more.
\end{tcolorbox}
\section[\texorpdfstring{\ta{ஒத்திவை} (ottivai)}{ஒத்திவை (ottivai)}]{\ta{ஒத்திவை} (ottivai) — to postpone; to put off}

\label{lesson:TA-C400-ottivai}



\begin{quote}
Before the new one: say the Tamil for a railway track, then the Tamil for arrival.
\end{quote}

\subsection*{You'll want to know: \ta{ஒத்திவை}}
\textbf{\ta{ஒத்திவை}} — \emph{ottivai} — \textquotedblleft{}to postpone; to put off\textquotedblright{}.

\begin{quote}
\ta{மழை} \ta{பெய்ததால்} \ta{நாங்கள்} \ta{பயணத்தை} \ta{ஒத்திவைத்தோம்}.
\end{quote}

\begin{quote}
\emph{maḻai peytatāl nāṅkaḷ payaṇattai ottivaittōm.}
\end{quote}

\begin{quote}
Because it rained, we put off the trip.
\end{quote}

Past (I …): \ta{ஒத்திவைத்தேன்}, \emph{ottivaittēṉ}.

What is put off takes -\ta{ஐ}: \ta{பயணத்தை} \ta{ஒத்திவை}, to postpone a journey.

\begin{cousinweb}
\textbf{\ta{ஒத்திவை}} (\emph{ottivai}) is built from \emph{otti}, the joining form of \ta{ஒத்து}, \textquotedblleft{}to move aside, make room\textquotedblright{} (DEDR 973), and \ta{வை}, \textquotedblleft{}to put\textquotedblright{}: to put a thing aside for later.
\end{cousinweb}

\begin{grammarlens}[title={What you've built}]
One more word for getting around.
\end{grammarlens}

\subsection*{Your turn}
\begin{itemize}
\raggedright
  \item \emph{Say it:} \emph{ottivai}
  \item \emph{Say it:} \emph{ottivai}, once more
  \item \emph{Say it:} \emph{maḻai peytatāl nāṅkaḷ payaṇattai ottivaittōm.} — Because it rained, we put off the trip.
  \item \emph{Recall:} say the Tamil for a railway track, then the Tamil for arrival, then say \emph{ottivai} again
\end{itemize}

\begin{tcolorbox}[breakable,colback=teal!4,colframe=teal!35!black,title={Before you move on}]
What does \ta{ஒத்திவை} mean? (\textquotedblleft{}To postpone; to put off\textquotedblright{}.) Say it once more.
\end{tcolorbox}
\section[\texorpdfstring{\ta{மாறு} (māṟu)}{மாறு (māṟu)}]{\ta{மாறு} (māṟu) — to switch; to change}

\label{lesson:TA-C400-maru-change}



\begin{quote}
Before the new one: say the Tamil for arrival, then the Tamil for to postpone.
\end{quote}

\subsection*{You'll want to know: \ta{மாறு}}
\textbf{\ta{மாறு}} — \emph{māṟu} — \textquotedblleft{}to switch; to change\textquotedblright{}.

\begin{quote}
\ta{அடுத்த} \ta{நிலையத்தில்} \ta{நாங்கள்} \ta{ரயில்} \ta{மாற} \ta{வேண்டும்}.
\end{quote}

\begin{quote}
\emph{aṭutta nilaiyattil nāṅkaḷ rayil māṟa vēṇṭum.}
\end{quote}

\begin{quote}
We have to change trains at the next station.
\end{quote}

Past (I …): \ta{மாறினேன்}, \emph{māṟiṉēṉ}.

\ta{மாறு} is to change or become different oneself; \ta{மாற்று} is to change something.

\begin{cousinweb}
\textbf{\ta{மாறு}} (\emph{māṟu}) is native Dravidian (DEDR 4834); the same article holds the transitive \ta{மாற்று}, \textquotedblleft{}to change something\textquotedblright{}.
\end{cousinweb}

\begin{grammarlens}[title={What you've built}]
One more word for getting around.
\end{grammarlens}

\subsection*{Your turn}
\begin{itemize}
\raggedright
  \item \emph{Say it:} \emph{māṟu}
  \item \emph{Say it:} \emph{māṟu}, once more
  \item \emph{Say it:} \emph{aṭutta nilaiyattil nāṅkaḷ rayil māṟa vēṇṭum.} — We have to change trains at the next station.
  \item \emph{Recall:} say the Tamil for arrival, then the Tamil for to postpone, then say \emph{māṟu} again
\end{itemize}

\begin{tcolorbox}[breakable,colback=teal!4,colframe=teal!35!black,title={Before you move on}]
What does \ta{மாறு} mean? (\textquotedblleft{}To switch; to change\textquotedblright{}.) Say it once more.
\end{tcolorbox}
